\exercisesection{Complements of Linear Algebra}

\begin{enumerate}
\def\labelenumi{\arabic{enumi})}
\tightlist
\item
  Let
\end{enumerate}

\[
\mathbf {M} \rightarrow \mathbf {N} \xrightarrow {t} \mathbf {P} \rightarrow \mathbf {Q}
\]

\[
\begin{array}{c c c c c} \downarrow u & & \downarrow v & & \downarrow w \\ \downarrow & & \downarrow & & \downarrow s \end{array}
\]

\[
\mathrm{M} ^ {\prime} \rightarrow \mathrm{N} ^ {\prime} \xrightarrow {t ^ {\prime}} \mathrm{P} ^ {\prime} \rightarrow \mathrm{Q} ^ {\prime}
\]

a commutative diagram of A-modules where the rows are exact; suppose that \(u\) is surjective and \(s\) injective. Show that:

\[
\operatorname{Ker} (w) = t (\operatorname{Ker} (v)) \quad \operatorname{Im} (v) = t ^ {\prime - 1} (\operatorname{Im} (w)).
\]

\begin{enumerate}
\def\labelenumi{\arabic{enumi})}
\setcounter{enumi}{1}
\tightlist
\item
  Let \(u: \mathbf{M} \to \mathbf{N}, v: \mathbf{N} \to \mathbf{P}\) two homomorphisms of A-modules. Show that there exists an exact sequence:
\end{enumerate}

\[
0 \rightarrow \operatorname{Ker} (u) \rightarrow \operatorname{Ker} (v \circ u) \rightarrow \operatorname{Ker} (v) \rightarrow \operatorname{Coker} (u) \rightarrow \operatorname{Coker} (v \circ u) \rightarrow \operatorname{Coker} v \rightarrow 0.
\]

\begin{enumerate}
\def\labelenumi{\arabic{enumi})}
\setcounter{enumi}{2}
\tightlist
\item
  In the diagram of A-modules:
\end{enumerate}

\[
\begin{array}{c} \mathrm{P} \xleftarrow {} \mathrm{M} \\ \Big \downarrow \quad \mathrm{R} ^ {\swarrow} \Big \downarrow \\ \mathrm{N} \xleftarrow {} \mathrm{Q} \end{array}
\]

Does the commutativity of the two triangles extracted from the diagram necessarily imply that of the diagram?

\begin{enumerate}
\def\labelenumi{\arabic{enumi})}
\setcounter{enumi}{3}
\item
  Let \(k\) a field, \(B = k[T]\) , let A be the subring \(k[T^2, T^3] \subset k[T]\) . Show that the A-module \(B\) is torsion-free, but is not a flat A-module.
\item
  Let \(k\) a field, \(\mathbf{C} = k[\mathbf{X}, \mathbf{Y}]\) ; consider the subrings \(\mathbf{A} = k[\mathbf{X}^4, \mathbf{Y}^4]\) , \(\mathbf{B} = k[\mathbf{X}^4, \mathbf{X}^3\mathbf{Y}, \mathbf{XY}^3, \mathbf{Y}^4]\) and \(\tilde{\mathbf{B}} = k[\mathbf{X}^4, \mathbf{X}^3\mathbf{Y}, \mathbf{X}^2\mathbf{Y}^2, \mathbf{XY}^3, \mathbf{Y}^4]\) of \(\mathbf{C}\) . The inclusions \(\mathbf{A} \subset \mathbf{B} \subset \tilde{\mathbf{B}}\) make it possible to define on \(\mathbf{B}\) and \(\tilde{\mathbf{B}}\) A-module structures. Show that \(\tilde{\mathbf{B}}\) is a flat A-module, while the A-module \(\mathbf{B}\) is not flat.
\item
  Show that an A-module E is flat if and only if it satisfies the following condition (cf. X, p. 15, remark):
\end{enumerate}

(ii'') For every finite family \((c_{i})_{i\in I}\) of elements of A, every solution \(e = (e_i)_{i\in I}\) of the equation \(\sum_{i\in I}c_i e_i = 0\) can be written \(b_1 z_1 + \cdots + b_n z_n\) , where \(b_1, \ldots, b_n \in E\) and where, for \(r = 1, \ldots, n\) , \(z_r = (z_{r,i})_{i \in I}\) is a solution of the equation \(\sum_{i=1}^{n} z_{r,i} c_i = 0\) .

\begin{enumerate}
\def\labelenumi{\arabic{enumi})}
\setcounter{enumi}{6}
\tightlist
\item
  Let E be a monoid and S a subsemigroup of E that is cancellative and satisfies the hypotheses of Exercise 17 of III, p. 121, so that there exists a “monoid of fractions” \(\overline{\mathbf{E}}\) and an injective homomorphism \(\varepsilon: \mathbf{E} \to \overline{\mathbf{E}}\) . Let A be a commutative ring. From \(\varepsilon\) one deduces a homomorphism of monoid algebras \(u: \mathbf{A}^{(\mathbf{E})} \to \mathbf{A}^{(\overline{\mathbf{E}})}\) which allows one to regard \(\mathbf{A}^{\overline{\mathbf{E}}}\) as a \(\mathbf{A}^{(\mathbf{E})}\) -module on the left. Show that this module is flat over \(\mathbf{A}^{(\mathbf{E})}\) (write \(\mathbf{A}^{\overline{\mathbf{E}}}\) as a filtered inductive limit of free \(\mathbf{A}^{(\mathbf{E})}\) -modules of rank 1).
\end{enumerate}

* 8) Let A be the ring of continuous real-valued functions on the interval {[}0, 1{]}.

\begin{enumerate}
\def\labelenumi{\alph{enumi})}
\item
  Show that the ideal of functions vanishing at 0 is a flat A-module (show that it is a filtered inductive limit of the ideals Af, where f is a function which vanishes only at 0).
\item
  Show that the ideal I of functions vanishing in a neighborhood of 0 is a projective A-module (show that there exists a sequence \((f_n)_{n\geqslant 1}\) , \(f_n\in \mathbf{I}\) , such that \(\operatorname {Supp}(f_n)\subset \left[\frac{1}{n + 2},\frac{1}{n}\right]\) and \(\sum_{n}f_{n}(x) = 1\) for \(x\neq 0\) ; deduce
\end{enumerate}

that every element \(h\) of I can be written \(h = \sum_{n}^{t_{h}^{\prime}} (f_n h) e_n\) with \(e_n = \sum_{r \leqslant n+1}^{n} f_r\) , and use II, p. 46, prop. 12 . *

\begin{enumerate}
\def\labelenumi{\arabic{enumi})}
\setcounter{enumi}{8}
\tightlist
\item
  Let P be an A-module. Show that the following conditions are equivalent:
\end{enumerate}

α) For a sequence \(M' \xrightarrow{u} M \xrightarrow{v} M''\) of right A-modules to be exact, it is necessary and sufficient that the sequence

\[
\mathbf {M} ^ {\prime} \otimes_ {\mathbf {A}} \mathbf {P} \xrightarrow {u \otimes 1} \mathbf {M} \otimes_ {\mathbf {A}} \mathbf {P} \xrightarrow {v \otimes 1} \mathbf {M} ^ {\prime \prime} \otimes_ {\mathbf {A}} \mathbf {P}
\]

be exact.

\(\beta)\) P is flat, and for every nonzero right A-module M, one has M \(\otimes_{A} P \neq (0)\) .

γ) P is flat, and for all right A-modules M, N, the canonical homomorphism

\[
\operatorname{Hom} _ {\mathbf {A}} (\mathbf {M}, \mathbf {N}) \rightarrow \operatorname{Hom} _ {\mathbf {Z}} (\mathbf {M} \otimes_ {\mathbf {A}} \mathbf {P}, \mathbf {N} \otimes_ {\mathbf {A}} \mathbf {P})
\]

is injective.

δ) P is flat, and for every maximal right ideal m of A, one has mP ≠ P.

One then says that P is a faithfully flat A-module.

\begin{enumerate}
\def\labelenumi{\arabic{enumi})}
\setcounter{enumi}{9}
\tightlist
\item
  \begin{enumerate}
  \def\labelenumii{\alph{enumii})}
  \tightlist
  \item
    Show that a faithfully flat module is faithful (II, p. 28).
  \end{enumerate}
\end{enumerate}

\begin{enumerate}
\def\labelenumi{\alph{enumi})}
\setcounter{enumi}{1}
\item
  The Z-module Q is faithful and flat, but not faithfully flat.
\item
  Let \((p_n)\) the strictly increasing sequence of prime numbers, and let A be the product ring \(\prod_{n} \mathbf{Z}/p_n \mathbf{Z}\) .
\end{enumerate}

Show that the direct sum E of the \(Z/p_{n}Z\) is a faithful projective A-module which is not faithfully flat (observe that E is an ideal of A such that \(E^{2}=E\) ).

\begin{enumerate}
\def\labelenumi{\arabic{enumi})}
\setcounter{enumi}{10}
\item
  Let \(0 \to \mathbf{P}' \to \mathbf{P} \to \mathbf{P}'' \to 0\) an exact sequence of A-modules. Suppose that \(\mathbf{P}'\) and \(\mathbf{P}''\) are flat, and that one of them is faithfully flat. Show that \(\mathbf{P}\) is faithfully flat.
\item
  Let \(\rho: A \to B\) a ring homomorphism that makes \(B\) an \(A\) faithfully flat right module; let \(M\) be an \(A\) module.
\end{enumerate}

\begin{enumerate}
\def\labelenumi{\alph{enumi})}
\item
  Show that for M to be finitely generated (resp. of finite presentation), it is necessary and sufficient that the same be true of the B-module \(\mathbf{B} \otimes_{\mathbf{A}} \mathbf{M}\) .
\item
  We assume that \(\rho(A)\) is contained in the center of B. Show that M is flat (resp. faithfully flat, resp. finitely generated projective) if and only if the B-module \(B \otimes_{A} M\) is.
\end{enumerate}

\begin{enumerate}
\def\labelenumi{\arabic{enumi})}
\setcounter{enumi}{12}
\tightlist
\item
  Show that every finitely generated flat module P over an integral domain A is projective (consider a presentation \(0 \rightarrow R \xrightarrow{L} L \xrightarrow{P} 0\) , where L is finitely generated free; choose a finitely generated submodule \(R'\) of R such that \(R/R'\) be torsion, and apply Theorem 1, p. 14).
\end{enumerate}

¶ 14) a) Let \(0 \to \mathbf{R} \to \mathbf{L} \to \mathbf{E} \to 0\) an exact sequence of A-modules, where L is a free A-module; let \((e_{a})\) a basis of L. Show that the following conditions are equivalent:

β) For every \(x \in \mathbb{R}\) , if \(\mathfrak{a}_{x}\) is the right ideal generated by the components of \(x\) on the basis \((e_{\alpha})\) , we have \(x \in \mathfrak{a}_{x} \mathbb{R}\) .

\(\gamma)\) For every \(x \in \mathbb{R}\) , there exists a homomorphism \(u_x: \mathbb{L} \to \mathbb{R}\) such that \(u_x(x) = x\) .

δ) For every finite sequence \((x_{i})_{1\leqslant i\leqslant n}\) of elements of R, there exists a homomorphism \(u:L\to R\) such that \(u(x_{i})=x_{i}\) for \(1\leqslant i\leqslant n\) (Use Thm. 1, p. 14.)

\begin{enumerate}
\def\labelenumi{\alph{enumi})}
\setcounter{enumi}{1}
\item
  Let r be the radical of A, and let \(0 \rightarrow R \rightarrow L \rightarrow E \rightarrow 0\) an exact sequence of A-modules, such that L is free. Suppose that E is flat and that R is contained in rL. Show that then R = 0 (with the notations of a), observe that \(a_{x}\) is a finitely generated ideal and that we have \(a_{x} = a_{x} r\) ).
\item
  Let E be a flat finitely generated A-module; suppose there exists a two-sided ideal b of A contained in the radical of A, such that E/bE is a free (A/b)-module. Show that E is then a free A-module (observe that there exists a finitely generated free A-module L such that L/bL is isomorphic to E/bE, and apply b)). In particular, every finitely generated flat module over a local ring is free.
\end{enumerate}

\begin{enumerate}
\def\labelenumi{\arabic{enumi})}
\setcounter{enumi}{14}
\tightlist
\item
  Let A be a ring, \(a\) an element of A. Show that the following properties are equivalent: \(\alpha)\) \(a\in aAa\) .
\end{enumerate}

β) Aa is a direct summand in the module \(A_{s}\) .

γ) \(A_{s}/Aa\) is a flat A-module.

δ) For every right ideal b of A, we have b ∩ Aa = ba.

\begin{enumerate}
\def\labelenumi{\arabic{enumi})}
\setcounter{enumi}{15}
\tightlist
\item
  Let A be a ring. Show that the following properties are equivalent:
\end{enumerate}

α) Every element \(a \in A\) satisfies the equivalent properties of Exercise 15.

β) Every finitely generated left ideal is a direct summand of \(A_{s}\) .

γ) Every left A-module is flat.

δ) Every right A-module is flat.

We then say that A is an absolutely flat ring (cf. VIII, § 8, Exercise 15).

\begin{enumerate}
\def\labelenumi{\arabic{enumi})}
\setcounter{enumi}{16}
\tightlist
\item
  Let A be an absolutely flat ring (Exercise 16).
\end{enumerate}

\begin{enumerate}
\def\labelenumi{\alph{enumi})}
\item
  Let P be a projective A-module. Show that every finitely generated submodule of P is a direct summand of P (reduce to the case where P is finitely generated free).
\item
  Show that every projective A-module is a direct sum of cyclic submodules, isomorphic to cyclic ideals of A. (Use Kaplansky's theorem (II, p. 183, exercise 2) to reduce to the case where P admits a countable family of generators, then use a).)
\item
  Show that every ideal of A generated by a countable family of generators is projective.
\item
  Give an example of an absolutely flat ring A and a finitely generated A-module that is not projective.
\end{enumerate}

18) Let M be an A-module. Show that the following conditions are equivalent:

α) M is finitely generated (resp. finitely presented).

β) For every inductive system \((\mathbf{N}_{i}, u_{ji})\) of A-modules relative to a filtered preordered set I, the canonical homomorphism \(\varinjlim\operatorname{Hom}_{\mathbf{A}}(\mathbf{M}, \mathbf{N}_{i}) \to \operatorname{Hom}_{\mathbf{A}}(\mathbf{M}, \varinjlim \mathbf{N}_{i})\) is injective (resp. bijective).

γ) For every family \((\mathbf{P}_{j})_{j\in J}\) of right A-modules, the canonical homomorphism

\[
\left(\prod_ {j \in J} P _ {j}\right) \otimes_ {A} M \rightarrow \prod_ {j \in J} \left(P _ {j} \otimes_ {A} M\right)
\]

(II, p. 61) is surjective (resp. bijective).

δ) For every set I, the canonical homomorphism \(A_{d}^{1} \otimes_{A} M \to M^{1}\) is surjective (resp. bijective). 19) Let A and B be two rings, \(\varphi : A \to B\) a homomorphism, M a B-module. We equip \(B \otimes_{A} M\) and \(\operatorname{Hom}_{A}(B, M)\) with left B-module structures induced from the B-module structure of B. One says that M is projective relative to A, or \(\varphi\) -projective, if the canonical surjection \(B \otimes_{A} M \to M\) admits a B-linear section; one says that M is injective relative to A, or \(\varphi\) -injective, if the canonical injection

\[
\mathbf {M} \rightarrow \operatorname{Hom} _ {\mathbf {A}} (\mathbf {B}, \mathbf {M})
\]

admits a B-linear retraction.

\begin{enumerate}
\def\labelenumi{\alph{enumi})}
\item
  If M is projective (resp. injective) as an A-module and \(\varphi\) -projective (resp. \(\varphi\) -injective), then M is a projective (resp. injective) B-module.
\item
  For every A-module N, the B-module \(B \otimes_{A} N\) (resp. \(\operatorname{Hom}_{A}(B, N)\) ) is \(\varphi\) -projective (resp. \(\varphi\) -injective).
\item
  We assume that \(\varphi(A)\) is contained in the center of B. Let N be an A-module, P a B-module \(\varphi\) -projective. Show that the B-module \(P \otimes_{A} N\) (resp. the right B-module \(\operatorname{Hom}_{A}(P, N)\) ) is \(\varphi\) -projective (resp. \(\varphi\) -injective).
\end{enumerate}

\begin{enumerate}
\def\labelenumi{\arabic{enumi})}
\setcounter{enumi}{19}
\item
  Let A be a principal ideal ring, \(a\) a nonzero element of A, B = A/Aa. Show that the B-module B \(_{s}\) is injective.
\item
  Show that the following conditions are equivalent:
\end{enumerate}

α) A is Noetherian.

β) Every A-module that is an inductive limit of injective A-modules is injective.

γ) Every A-module that is a direct sum of injective A-modules is injective.

(To prove that \(\gamma\) ) implies \(\alpha\) ), let \((\mathfrak{a}_n)_{n\geqslant 1}\) an increasing sequence of left ideals of A, \(\mathfrak{a} = \bigcup_{n}\mathfrak{a}_{n}\) ,

\(I_{n}\) an injective envelope of \(\mathfrak{a}/\mathfrak{a}_{n}\) . The natural homomorphism from \(\mathfrak{a}\) to \(\oplus I_{n}\) is the restriction of a homomorphism \(f: A \to \oplus I_{n}\) ; consider \(f(1)\) .

\begin{enumerate}
\def\labelenumi{\arabic{enumi})}
\setcounter{enumi}{21}
\tightlist
\item
  Let A be a principal ideal ring, K its field of fractions, \(\pi\) an extremal element of A. Show that the \(\pi\) -primary component of the A-module K/A is an injective envelope of the A-module A/ \(\pi^n\) A for every \(n \geqslant 1\) . ¶ 23) Let I be an injective A-module, E the ring End \(_A\) (I), r the set of elements \(u\) of E such that I is an injective envelope of Ker( \(u\) ).
\end{enumerate}

\begin{enumerate}
\def\labelenumi{\alph{enumi})}
\item
  Show that r is a two-sided ideal and that the ring E/r is absolutely flat (exercise 16). (Let u be an element of E, N a submodule of I such that N ∩ Ker u = 0 and maximal for this property; show that there exists v ∈ E such that vu(x) = x for all x ∈ N, and deduce that uvu - u ∈ r.)
\item
  Show that r is the radical of E (the inclusion r(E) ⊂ r follows from a); show that for every u ∈ r, 1 - u is bijective).
\item
  Assume that I is a direct sum of a finite family of indecomposable injectives. Show that the ring E/r is semisimple.
\end{enumerate}

Let \(i: \mathbf{A} \to \mathbf{I}\) an injective envelope of the A-module \(\mathbf{A}_{s}\) . We set \(\mathbf{E} = \operatorname{End}_{\mathbf{A}}(\mathbf{I})\) and \(\mathbf{B} = \operatorname{End}_{\mathbf{E}}(\mathbf{I})\) , so that \(\mathbf{B}\) is the bicommutant of the A-module I (VIII, § 3, def. 1). We consider \(\mathbf{A}\) as a subring of \(\mathbf{B}\) . One denotes \(j: \mathbf{B} \to \mathbf{I}\) and \(p: \mathbf{E} \to \mathbf{I}\) the maps defined by \(j(b) = b(i(1))\) and \(p(u) = u(i(1))\) for \(b \in \mathbf{B}\) , \(u \in \mathbf{E}\) .

\begin{enumerate}
\def\labelenumi{\alph{enumi})}
\item
  Show that \(p\) is a surjection and \(j\) an injection.
\item
  Show that the following conditions are equivalent:
\end{enumerate}

α) The map p is bijective.

β) The map j is bijective.

γ) The A-module B is injective.

If these conditions are satisfied, the mapping \(j^{-1} \circ p\) is an isomorphism from the ring E onto the opposite ring of B.

\begin{enumerate}
\def\labelenumi{\alph{enumi})}
\setcounter{enumi}{2}
\item
  Let \(\mathfrak{F}\) the set of left ideals \(a\) of \(A\) such that \(a \cap b \neq (0)\) for every nonzero left ideal \(b\) . One denotes \(s(A)\) the set of elements \(a\) of \(A\) for which there exists \(a \in \mathfrak{F}\) such that \(aa = 0\) . Show that \(s(A)\) is a two-sided ideal of \(A\) .
\item
  Show that the following conditions are equivalent:
\end{enumerate}

α) \(s(\mathbf{A}) = 0.\)

β) The ring E is semiprimitive.

γ) The ring B is absolutely flat.

If these conditions are satisfied, then the equivalent properties of \(b\) ) are satisfied.

(Let \(s(I)\) the set of elements \(x\) of I for which there exists \(\mathfrak{a} \in \mathfrak{F}\) such that \(\mathfrak{a} x = 0\) . Using Exercise 23, prove that \(p^{-1}(s(I))\) is equal to the radical \(r\) of E, whence the implication \(\beta) \Rightarrow \alpha\) ). Then show that \(\alpha\) implies \(r = \text{Ker } (p) = 0\) , which entails \(\beta\) and \(\gamma\) (in view of Exercise 23) as well as condition \(\alpha\) of \(b\) ). Finally show that \(\gamma\) implies \(\alpha\) .)

\begin{enumerate}
\def\labelenumi{\alph{enumi})}
\setcounter{enumi}{4}
\item
  If A is left Noetherian, the ideal s(A) is nilpotent (first prove that every element a of s(A) is nilpotent, by considering the left annihilators of the \(a^{n}\) ; then use Exercise 8, b) of VIII, § 9). ¶ 25) Assume that the ring A is left Noetherian and contains no nonzero nilpotent two-sided ideals. a) Prove that the ring B defined in Exercise 24 is semisimple (“Goldie's theorem”; use Exercise 23, c) and Exercise 24, d) and e)). If moreover (0) is a prime two-sided ideal in A (VIII, § 8, Exercise 19), show that B is simple.
\item
  One identifies A with a subring of B. Prove the following properties:
\end{enumerate}

\begin{enumerate}
\def\labelenumi{(\roman{enumi})}
\item
  Every element of A that is left cancellable is invertible in B (hence cancellable in A).
\item
  Every element \(b\) of \(\mathbf{B}\) can be written \(s^{-1}a\) , with \(s \in \mathbf{A}\) , \(a \in \mathbf{A}\) and \(s\) cancellable.
\end{enumerate}

One sometimes says that B is the total ring of left fractions of A.

(To prove (ii), reduce to the case where B is simple. Show then that the ideal (0) of A is prime. Let a be an ideal of F, a an element of a for which the right ideal aB is maximal among the ideals a'B (a' ∈ a), and x an element of B such that axa = a. Prove that B(1 - xa) ⊂ aB, whence

\[
(\mathbf {B} (1 - a x) \cap \mathfrak {a}) \cdot (\mathbf {B} (1 - x a) \cap \mathfrak {a}) = (0);
\]

deduce that \(a\) is cancellable. Conclude by taking for a the set of all \(a \in A\) such that \(ab \in A\) .) c) Assume that every nonzero element of \(A\) is cancellable. Show that \(B\) is a field, which is a left field of fractions for \(A\) (cf. I, p. 155, Exercise 15).

¶ 26) a) Show that the following conditions are equivalent:

α) The ring A is right Noetherian, and the A-module \(A_{s}\) is injective.

\(\alpha^{0}\) A is left Noetherian and the A-module \(A_{d}\) is injective.

β) The ring A is self-injective (VIII, § 2, Exercise 11).

(Prove that \(\alpha\) ) implies that every right ideal r satisfies \(r = r^0\) . Then show that \(r(A)\) is nilpotent, and deduce that A is right Artinian using Exercise 23. Next prove that the dual of a simple left A-module is simple. Deduce that A is left Noetherian and finally self-injective. b) Show that the ring A is self-injective if and only if every injective A-module is projective and every projective A-module is injective.

¶ 27) Assume the ring A is left Noetherian.

\begin{enumerate}
\def\labelenumi{\alph{enumi})}
\item
  Let I be an indecomposable injective A-module. Show that the set of annihilator ideals of a nonzero submodule of I admits a greatest element \(\mathfrak{A}(I)\) , which is a prime two-sided ideal (VIII, § 8, Exercise 19).
\item
  Show that for every prime two-sided ideal p of A, there exists an indecomposable injective A-module I such that \(\mathfrak{A}(I) = \mathfrak{p}\) (consider an indecomposable direct factor of the injective envelope of A/p).
\item
  Assume A is commutative. Show that every indecomposable injective A-module I is the injective envelope of A/U(I). Deduce that the map I ↦ U(I) defines a bijection from the set S of isomorphism classes of indecomposable injective A-modules onto the set of prime ideals of A. The inverse bijection associates to the ideal p the injective envelope of A/p.
\item
  Show that for the map I ↦ \(\mathfrak{A}(I)\) from \(\mathfrak{A}\) onto the set of prime two-sided ideals of A to be bijective, it suffices that A satisfy the following condition (H):
\end{enumerate}

\begin{enumerate}
\def\labelenumi{(\Alph{enumi})}
\setcounter{enumi}{7}
\tightlist
\item
  For every left ideal a of A, there exist elements \(x_{1}, \ldots, x_{k}\) of A/a such that
\end{enumerate}

\[
\operatorname{Ann} (\mathbf {A} / \mathfrak {a}) = \bigcap_ {i} \operatorname{Ann} (x _ {i}).
\]

\begin{enumerate}
\def\labelenumi{\alph{enumi})}
\setcounter{enumi}{4}
\tightlist
\item
  Show that condition (H) is satisfied in the following cases:
\end{enumerate}

α) All left ideals of A are two-sided.

β) A is Artinian.

γ) The center Z of A is a Noetherian ring and A is a finitely generated Z-module.

\begin{enumerate}
\def\labelenumi{\alph{enumi})}
\setcounter{enumi}{5}
\tightlist
\item
  Let k be a field, \(Q^{+}\) the additive monoid of rational numbers \(\geqslant 0\) , A the algebra \(k^{(\mathbf{Q}^{+})}\) , a a nonzero element of A, I the injective envelope of A/Aa. Show that the A-module I is indecomposable, but contains no submodule isomorphic to A/p for p prime (note that the set of ideals of A is totally ordered and that A contains only two prime ideals).
\end{enumerate}

\begin{enumerate}
\def\labelenumi{\arabic{enumi})}
\setcounter{enumi}{27}
\tightlist
\item
  Let A be a Noetherian ring, M an A-module, I an injective envelope of M, a direct sum of a family \((I_{\alpha})_{\alpha \in J}\) of indecomposable submodules. We denote by Ass (M) the set of prime two-sided ideals \(\mathfrak{U}(I_{\alpha})\) , for \(\alpha \in J\) (exercise 27); if \(p \in \operatorname{Ass}(M)\) , we say that p is associated with M.
\end{enumerate}

\begin{enumerate}
\def\labelenumi{\alph{enumi})}
\item
  We have Ass (M) ≠ ∅ if M ≠ 0; if M is finitely generated, the set Ass (M) is finite.
\item
  Show that a prime two-sided ideal p is associated with M if and only if M contains a submodule N such that Ann (N') = p for every nonzero submodule N' of N. If A is commutative, p is associated with M if and only if M contains a submodule isomorphic to A/p.
\item
  Let M' be a submodule of M. Show that we have
\end{enumerate}

\[
\operatorname{Ass} \left(\mathbf {M} ^ {\prime}\right) \subset \operatorname{Ass} (\mathbf {M}) \subset \operatorname{Ass} \left(\mathbf {M} ^ {\prime}\right) \cup \operatorname{Ass} \left(\mathbf {M} / \mathbf {M} ^ {\prime}\right).
\]

\begin{enumerate}
\def\labelenumi{\alph{enumi})}
\setcounter{enumi}{3}
\tightlist
\item
  We set \(J_{p} = \sum_{\mathfrak{M}(I_{\alpha}) = p} I_{\alpha}\) and \(Q(p) = \left( \sum_{q \neq p} J_{q} \right) \cap M\) Show that we have \(\bigcap_{p \in \operatorname{Ass}(M)} Q(p) = 0\) and \(\bigcap_{p \in E} Q(p) \neq 0\) for every subset E of Ass(M) distinct from Ass(M). For every \(p \in \operatorname{Ass}(M)\) , we have
\end{enumerate}

Ass \((\mathbf{M} / \mathbf{Q}(\mathfrak{p})) = \{\mathfrak{p}\}\) .

\begin{enumerate}
\def\labelenumi{\alph{enumi})}
\setcounter{enumi}{4}
\tightlist
\item
  Let a be an element of the center of A. Show that for the homothety of ratio a to be injective in M, it is necessary and sufficient that a belong to none of the prime ideals associated with M.
\end{enumerate}

Assume that the ring A is commutative and Noetherian. Let a be an ideal of A.

\begin{enumerate}
\def\labelenumi{\alph{enumi})}
\tightlist
\item
  Let H be an A-module such that every element of H is annihilated by a power of a. Show that for H to be injective, it suffices that for every integer \(n \geqslant 1\) every finitely generated A-module M annihilated by \(a^{n}\) and every submodule \(M'\) of M, the natural homomorphism \(\operatorname{Hom}_{\mathbf{A}}(\mathbf{M}, \mathbf{H}) \to \operatorname{Hom}_{\mathbf{A}}(\mathbf{M}', \mathbf{H})\) be surjective.
\end{enumerate}

(Use the fact that for every ideal b of A, there exists an integer \(k \geqslant 0\) such that \(\mathfrak{b} \cap \mathfrak{a}^k \subset \mathfrak{ba}\) (AC, III, § 3, no. 1, cor. 2).)

\begin{enumerate}
\def\labelenumi{\alph{enumi})}
\setcounter{enumi}{1}
\item
  Let M be an injective A-module. For \(n \geqslant 1\) , we denote by \(\mathbf{H}_{n}\) the submodule of M consisting of the elements annihilated by \(\alpha^{n}\) ; set \(\mathbf{H} = \bigcup \mathbf{H}_{n}\) . Show that H is an injective A-module.
\item
  Show that every element of the injective envelope of A/a is annihilated by a power of a. * 30) Assume the ring A is local, commutative, and Noetherian; denote its maximal ideal by m, and k = A/m. Let I be an A-module such that every element of I is annihilated by a power of m. For every A-module M, denote by T(M) the A-module Hom\_A (M, I) and c\_M : M → T(T(M)) the homomorphism defined by
\end{enumerate}

\[
\left(c _ {\mathbf {M}} (m)\right) (u) = u (m) \quad \text { for } \quad m \in \mathbf {M}  , \quad u \in \mathrm{T} (\mathbf {M})  .
\]

Show that the following properties are equivalent:

α) For every finite-length A-module M, the A-module T(M) is of finite length and the homomorphism \(c_{M}: M \to T(T(M))\) is bijective.

β) For every A-module M of finite length, one has long T(M) = long M.

γ) The A-module I is injective, and the A-modules k and T(k) are isomorphic.

δ) The A-module I is isomorphic to the injective envelope of k.

When these conditions are satisfied, one sometimes says that I is a dualizing module for the ring A. (Show that property γ) is equivalent to each of the other three: to see that α) or β) implies γ), use Exercise 29.)

\begin{enumerate}
\def\labelenumi{\arabic{enumi})}
\setcounter{enumi}{30}
\item
  With the hypotheses of the previous exercise, let B be a local, commutative, Noetherian ring, and let \(\rho: A \to B\) a ring homomorphism that makes B a finitely generated A-module. If I is a dualizing module for A, show that the B-module \(\mathrm{Hom}_{\mathbf{A}}(\mathbf{B}, I)\) is dualizing for B. In particular, for every ideal \(a\) of A, the submodule of I consisting of elements annihilated by \(a\) is a dualizing module for A/a.
\item
  \begin{enumerate}
  \def\labelenumii{\alph{enumii})}
  \tightlist
  \item
    We keep the hypotheses of Exercise 30; we further assume that A contains a field \(k_0\) such that \(k\) is a finite-dimensional \(k_0\) -vector space. Show that the A-module \(l = \varinjlim_{n} \operatorname{Hom}_{k_0}(A/m^n, k_0)\)
  \end{enumerate}
\end{enumerate}

is dualizing for A.

\begin{enumerate}
\def\labelenumi{\alph{enumi})}
\setcounter{enumi}{1}
\tightlist
\item
  Let \(k\) a field of characteristic zero, B the ring of formal power series \(k[[\mathrm{T}_{1}, \ldots, \mathrm{T}_{n}]]\) . We endow the group \(\mathbf{I} = k[\mathrm{X}_{1}, \ldots, \mathrm{X}_{n}]\) with a B-module structure by setting \(\mathrm{T}_{i} \cdot \mathbf{P} = \partial \mathbf{P} / \partial \mathbf{X}_{i}\) for all \(\mathbf{P} \in k[\mathrm{X}_{1}, \ldots, \mathrm{X}_{n}]\) . Show that I is an injective envelope of \(k\) .
\end{enumerate}

